\documentclass[11pt]{article}
\usepackage[utf8]{inputenc}
\usepackage[T1]{fontenc}
\usepackage{amsmath,amssymb,amsthm}
\usepackage{graphicx}
\usepackage{hyperref}
\usepackage{booktabs}
\usepackage{algorithm}
\usepackage{algorithmic}
\usepackage{natbib}
\usepackage[margin=1in]{geometry}

\DeclareMathOperator{\eml}{eml}
\newcommand{\bx}{\mathbf{x}}
\newcommand{\bg}{\mathbf{g}}
\newcommand{\bK}{\mathbf{K}}
\newcommand{\bC}{\mathbf{C}}
\newcommand{\R}{\mathbb{R}}
\newcommand{\C}{\mathbb{C}}

\title{Koopman Operator Lifting via EML Observable Dictionaries:\\
A PySINDy-Integrated Framework with CTF Benchmarking}
\author{Jim Beaver}
\date{April 2026 --- Working Draft}

\begin{document}
\maketitle

\begin{abstract}
We introduce a framework for learning interpretable Koopman operator
representations using the EML (Exp-Minus-Log) Sheffer operator
$\eml(x,y) = \exp(x) - \ln(y)$ as the basis for observable dictionaries.
Since every elementary function can be expressed as a binary tree of identical
EML nodes, parameterized EML trees provide a \emph{complete}, gradient-trainable
search space for Koopman eigenfunctions.  After Gumbel-softmax annealing and
weight snapping, the learned observables are exact symbolic expressions ---
yielding closed-form Koopman eigenfunctions that black-box methods cannot produce.
We package EML observables as a PySINDy-compatible feature library and validate
against the dynamicsai.org Common Task Framework (CTF) 12-metric benchmark on the
Lorenz system, comparing to EDMD, deep Koopman autoencoders, and PySR baselines.
\end{abstract}

% ============================================================================
\section{Introduction}
\label{sec:intro}

% TODO: Expand with motivation, context, and contributions.

The Koopman operator provides a linear (but infinite-dimensional) description of
nonlinear dynamical systems.  A central challenge is selecting the dictionary of
observable functions that span a finite-dimensional invariant subspace.

Odrzywo{\l}ek (2026) discovered that the binary operator $\eml(x,y) = \exp(x) - \ln(y)$,
paired with the constant 1, generates \emph{every} standard elementary function ---
the continuous-math analog of the NAND gate for Boolean logic.  This provides
a natural, complete basis for Koopman observable dictionaries.

\paragraph{Contributions.}
\begin{itemize}
    \item Parameterized EML trees as gradient-trainable Koopman observables with
          exact symbolic snap.
    \item GPU-optimized Taylor-series EML primitives (zero transcendental calls).
    \item PySINDy integration via \texttt{EMLLibrary}.
    \item CTF benchmark evaluation (E1--E12) on Lorenz.
\end{itemize}


% ============================================================================
\section{Background}
\label{sec:background}

\subsection{The EML Operator}

% TODO: Summarize Odrzywołek's result, key identities.

\subsection{Koopman Operator Theory}

% TODO: Standard Koopman setup, EDMD, deep Koopman.

\subsection{SINDy and PySINDy}

% TODO: Brief overview of sparse identification.


% ============================================================================
\section{Method}
\label{sec:method}

\subsection{EML Trees as Koopman Observables}

% TODO: Master formula, Gumbel-softmax routing, three-phase training.

\subsection{Taylor-Series GPU Primitives}

% TODO: Horner-form exp/ln, range reduction, custom autograd.

\subsection{PySINDy Integration}

% TODO: EMLLibrary API, enumerative vs pre-trained modes.


% ============================================================================
\section{Experimental Setup}
\label{sec:experiments}

\subsection{CTF Benchmark Framework}

% TODO: Describe CTF, E1-E12 metrics, Lorenz dataset.

\subsection{Baselines}

% TODO: EDMD-Poly, EDMD-RBF, Deep Koopman, PySR.


% ============================================================================
\section{Results}
\label{sec:results}

% TODO: Tables, radar plots, prediction rollout figures, formula examples.

\begin{table}[h]
\centering
\caption{CTF Lorenz E1 (short-term forecasting) comparison.}
\label{tab:e1}
\begin{tabular}{lcc}
\toprule
Method & E1 Score & Parameters \\
\midrule
EML-Koopman (ours) & --- & --- \\
EDMD-Poly & --- & --- \\
EDMD-RBF & --- & --- \\
Deep Koopman & --- & --- \\
PySR & --- & --- \\
\midrule
LSTM (CTF published) & 64.54 & --- \\
DeepONet (CTF published) & 57.80 & --- \\
Reservoir (CTF published) & 54.87 & --- \\
\bottomrule
\end{tabular}
\end{table}


% ============================================================================
\section{Discussion}
\label{sec:discussion}

% TODO: Interpretability advantages, limitations, future work.


% ============================================================================
\section{Conclusion}
\label{sec:conclusion}

% TODO: Summary of contributions and outlook.


\bibliographystyle{plainnat}
% \bibliography{references}  % uncomment when .bib file is ready

\end{document}
